% !TEX program = pdflatex
% ============================================================
%  Robust Neural-Network Classification: Complete Theory,
%  Mathematical Derivations, and Novel Contributions
%  Date: 12 April 2026
%  Authors: [add names]
% ============================================================
\documentclass[11pt,a4paper]{article}

\usepackage[margin=1in]{geometry}
\usepackage{amsmath,amssymb,amsthm,mathtools}
\usepackage{booktabs,multirow,array}
\usepackage{graphicx}
\usepackage{hyperref}
\usepackage{cleveref}
\usepackage{microtype}
\usepackage{enumitem}
\usepackage{xcolor}
\usepackage{algorithm}
\usepackage{algpseudocode}
\usepackage{tcolorbox}
\usepackage{bbm}

\tcbuselibrary{theorems,skins}

\hypersetup{colorlinks=true,linkcolor=blue!70!black,citecolor=green!50!black,urlcolor=blue!50!black}

% ── Theorem environments ────────────────────────────────────────────────────
\newtheorem{theorem}{Theorem}[section]
\newtheorem{proposition}[theorem]{Proposition}
\newtheorem{lemma}[theorem]{Lemma}
\newtheorem{corollary}[theorem]{Corollary}
\newtheorem{definition}[theorem]{Definition}
\newtheorem{remark}[theorem]{Remark}

% ── Math macros ─────────────────────────────────────────────────────────────
\newcommand{\softmax}{\operatorname{softmax}}
\newcommand{\sign}{\operatorname{sign}}
\newcommand{\one}{\mathbf{1}}
\newcommand{\R}{\mathbb{R}}
\newcommand{\E}{\mathbb{E}}
\newcommand{\Prob}{\mathbb{P}}
\newcommand{\N}{\mathcal{N}}
\newcommand{\Lcal}{\mathcal{L}}
\newcommand{\norm}[1]{\left\|#1\right\|}
\newcommand{\abs}[1]{\left|#1\right|}
\newcommand{\ind}[1]{\mathbbm{1}\!\left[#1\right]}
\newcommand{\py}{p_y}
\newcommand{\pb}{\mathbf{p}}
\newcommand{\zb}{\mathbf{z}}
\newcommand{\xb}{\mathbf{x}}
\newcommand{\Tb}{\mathbf{T}}
\newcommand{\eb}{\mathbf{e}}
\newcommand{\half}{\tfrac{1}{2}}
% S-Divergence loss parameters
\newcommand{\Asd}{A_{\beta\lambda}}
\newcommand{\Bsd}{B_{\beta\lambda}}

\title{\textbf{Robust Neural Classification under Label Noise and Adversarial Perturbations:}\\
Theory, S-Divergence Unification, and Experimental Design}
\author{(12 April 2026)}
\date{}

\begin{document}
\maketitle

\begin{abstract}
This document provides a \textbf{self-contained, mathematically rigorous} treatment of robust loss functions for deep neural-network classification. We cover:
(1) a formal model of label noise and adversarial perturbations;
(2) complete derivations of all seven loss families used in our experiments—CCE, MAE, GCE, TDPDSCCE, SCE, SDIV, and Forward Correction;
(3) \textbf{a unification theorem} showing that the S-Divergence (SDIV) is a superfamily containing CCE, DPD (TDPDSCCE), and MAE as limiting or special cases;
(4) gradient-magnitude analysis explaining \emph{why} each loss confers robustness;
(5) an \textbf{influence-function} perspective quantifying the leverage of mislabeled samples;
(6) novel theoretical contributions: a \textbf{dual robustness bound} under simultaneous label noise and adversarial perturbation, and a curriculum-annealing proposal;
(7) experimental design guidelines for reproducible, conference-quality results.

This document is intended as the living mathematical companion to the code files
\texttt{12April2026\_RobustNN\_Experiments.py} and \texttt{12April2026\_RobustNN\_Research.ipynb}.
\end{abstract}

\tableofcontents
\newpage

%==============================================================================
\section{Notation and Probabilistic Setup}
\label{sec:notation}
%==============================================================================

Let $\mathcal{X} \subseteq \R^d$ be the input space and $\mathcal{Y} = \{1,\ldots,C\}$ the label set.
A neural network $f_\theta : \mathcal{X} \to \R^C$ maps inputs to \emph{logits} $\zb = f_\theta(\xb)$,
and the \emph{predicted probability vector} is
\[
  \pb = \softmax(\zb), \qquad p_k = \frac{e^{z_k}}{\sum_{j=1}^C e^{z_j}}, \quad k=1,\ldots,C.
\]
For a sample with true label $y \in \mathcal{Y}$ we write $\py = p_y$ for the predicted probability
on the correct class. The one-hot encoding is $\eb_y \in \{0,1\}^C$.

\paragraph{Clean-data risk.}
The population risk under a loss $\ell$ and distribution $P_{XY}$ is
$R(\theta) = \E[\ell(\pb, y)]$.

%==============================================================================
\section{Two Robustness Challenges: Label Noise and Adversarial Attacks}
\label{sec:challenges}
%==============================================================================

Understanding these two challenges is essential before studying which loss functions
address them, and how.

%------------------------------------------------------------------------------
\subsection{What is an Adversarial Attack?}
\label{sec:adversarial}
%------------------------------------------------------------------------------

An \textbf{adversarial example} is an input $\xb' = \xb + \boldsymbol{\delta}$ where the perturbation
$\boldsymbol{\delta}$ is \emph{imperceptibly small} (measured in $\ell_\infty$ or $\ell_2$ norm) yet
causes the model to misclassify. Unlike label noise (which corrupts training labels),
adversarial attacks corrupt \emph{test inputs}.

\begin{definition}[Adversarial Example]
Given a correctly classified input $\xb$, an $\epsilon$-adversarial example satisfies:
\[
  \xb' = \xb + \boldsymbol{\delta}, \quad \norm{\boldsymbol{\delta}}_\infty \leq \epsilon, \quad
  \arg\max_k p_k(f_\theta(\xb')) \neq y.
\]
\end{definition}

\paragraph{Why adversarial examples exist.}
Neural networks are locally linear in their decision boundaries (Goodfellow et al., 2014).
In high-dimensional input spaces, tiny steps in the gradient direction of the loss can
traverse large distances in function space.

\paragraph{Fast Gradient Sign Method (FGSM).}
The FGSM attack (Goodfellow et al., 2014) constructs $\boldsymbol{\delta}$ with a single gradient step:
\begin{equation}
  \label{eq:fgsm}
  \xb' = \xb + \epsilon \cdot \sign\!\left(\nabla_{\xb}\, \ell\big(\pb(\xb), y\big)\right).
\end{equation}
This maximizes the loss in the $\ell_\infty$-ball of radius $\epsilon$ around $\xb$.

\paragraph{Key insight for our experiments.}
FGSM requires computing $\nabla_{\xb} \ell$ — the gradient of the loss w.r.t.\ the \emph{input}.
Robust losses (GCE, SDIV) produce smaller input gradients at high confidence, potentially
reducing the effective attack strength even under the same $\epsilon$.
Our FGSM battery (Battery C, MNIST/CIFAR) directly tests this hypothesis.

\begin{tcolorbox}[colback=blue!5!white,colframe=blue!60!black,title=Intuition: Adversarial Attack]
\textbf{Analogy:} A perfectly drawn "7" that fools the model when you add a small
invisible $\epsilon$-size pattern overlaid on it. It looks identical to humans but the
model's local linear approximation sends it across a class boundary.
The key parameters are: $\epsilon$ (attack strength) and which loss was used to train the model.
\end{tcolorbox}

%------------------------------------------------------------------------------
\subsection{What is Label Noise?}
\label{sec:labelnoise}
%------------------------------------------------------------------------------

\textbf{Label noise} is corruption of training labels: $\tilde{y} \neq y^\star$.
Unlike adversarial perturbations, label noise affects training, not test inputs.

\begin{definition}[Noise Transition Matrix]
A row-stochastic matrix $\Tb \in [0,1]^{C\times C}$ with $\sum_j T_{ij} = 1$ describes
\emph{instance-independent} label noise:
\[
  T_{ij} = \Prob(\tilde{y} = j \mid y^\star = i).
\]
\end{definition}

\paragraph{Symmetric (uniform) noise.}
Each label is flipped to a uniformly random wrong class with probability $\eta$:
\[
  T_{ij} = \begin{cases} 1 - \eta & \text{if } i = j, \\ \eta/(C-1) & \text{if } i \neq j. \end{cases}
\]

\paragraph{Class-dependent (cyclic) noise.}
\[
  T_{ij} = \begin{cases} 1 - \eta & \text{if } j = i, \\ \eta & \text{if } j = (i+1)\bmod C, \\ 0 & \text{otherwise.} \end{cases}
\]

\paragraph{Why standard CCE fails.}
With probability $\eta$ the training label is wrong. CCE drives $\py \to 1$
for the (wrong) observed label, \emph{memorizing noise}. The phenomenon is
well-documented: deep networks can fit completely random labels in finite time
\cite{zhang2018gce}.

\begin{tcolorbox}[colback=orange!5!white,colframe=orange!60!black,title=Intuition: Label Noise vs Adversarial]
\textbf{Label noise} = wrong answers on the exam answer sheet. The teacher
(loss function) trains the student to memorize wrong answers. \\
\textbf{Adversarial attack} = the exam questions are subtly modified at test time to
trick the student, even though training was clean.
\end{tcolorbox}

%==============================================================================
\section{Loss Functions: Definitions, Derivations, and Intuitions}
\label{sec:losses}
%==============================================================================

All losses operate on the predicted softmax probabilities $\pb = \softmax(\zb)$ and
the observed label $y$. We write $p_y$ for the predicted mass on the observed class.

%------------------------------------------------------------------------------
\subsection{Categorical Cross-Entropy (CCE) — The Baseline}
\label{sec:cce}
%------------------------------------------------------------------------------

\paragraph{Definition.}
\begin{equation}
  \label{eq:cce}
  \Lcal_{\mathrm{CCE}}(\pb, y) = -\log \py.
\end{equation}

\paragraph{Statistical interpretation.}
CCE is the negative log-likelihood of a multinomial model:
minimizing $\E[\Lcal_{\mathrm{CCE}}]$ is equivalent to minimizing
the KL divergence $D_{\mathrm{KL}}(\eb_y \| \pb)$.

\paragraph{Gradient w.r.t.\ logits.}
\begin{equation}
  \label{eq:cce_grad}
  \frac{\partial \Lcal_{\mathrm{CCE}}}{\partial z_k}
  = p_k - \ind{k=y}.
\end{equation}
This is the celebrated ``softmax cross-entropy gradient'' — it is \textbf{unbounded}
as $p_y \to 0$ because $\partial \Lcal/\partial p_y = -1/p_y \to -\infty$.

\paragraph{Failure under noise.}
If a sample has noisy label $\tilde{y} \neq y^\star$, CCE applies full gradient force
to push $p_{\tilde y} \to 1$, with force $1/p_{\tilde y}$ that grows as the network
(correctly) places low mass on $\tilde y$. This is the memorization mechanism.

\paragraph{Scale.}
$\Lcal_{\mathrm{CCE}} \in [0, +\infty)$. At initialization with $\pb \approx (1/C,\ldots,1/C)$,
$\Lcal_{\mathrm{CCE}} \approx \log C$. For $C=10$ (CIFAR-10), $\log 10 \approx 2.3$.

%------------------------------------------------------------------------------
\subsection{Mean Absolute Error on Probabilities (MAE)}
\label{sec:mae}
%------------------------------------------------------------------------------

\paragraph{Definition.}
\begin{equation}
  \label{eq:mae}
  \Lcal_{\mathrm{MAE}}(\pb, y) = 1 - \py.
\end{equation}

\paragraph{Properties.}
$\Lcal_{\mathrm{MAE}} \in [0,1]$: \textbf{bounded}. The gradient $\partial \Lcal/\partial \py = -1$
is constant — independent of the current confidence $\py$. This means wrong labels
(low $p_y$ in a well-calibrated model) get the same gradient as right labels.

\paragraph{Note on scale.}
MAE is bounded in $[0,1]$. This is why in our plots MAE values are 0--1 whereas
CCE values are 0--2.3+. They must \emph{never} be compared on the same y-axis without normalization.

%------------------------------------------------------------------------------
\subsection{Generalized Cross-Entropy (GCE)}
\label{sec:gce}
%------------------------------------------------------------------------------

Proposed by Zhang and Sabuncu \cite{zhang2018gce} as a smooth interpolation
between CCE and MAE.

\paragraph{Definition.}
For $q \in (0, 1]$:
\begin{equation}
  \label{eq:gce}
  \Lcal_{\mathrm{GCE}}(\pb, y; q) = \frac{1 - \py^q}{q}.
\end{equation}

\paragraph{Limiting cases.}
\begin{align}
  q \to 0^+: \quad & \Lcal_{\mathrm{GCE}} = \frac{1-e^{q\log \py}}{q} \to -\log \py = \Lcal_{\mathrm{CCE}},\\
  q = 1: \quad      & \Lcal_{\mathrm{GCE}} = 1 - p_y = \Lcal_{\mathrm{MAE}}.
\end{align}
The first limit follows by L'Hôpital's rule: $\lim_{q\to 0}(1-e^{q\log p})/q = -\log p$.

\paragraph{Gradient magnitude (key robustness mechanism).}
\begin{equation}
  \frac{\partial \Lcal_{\mathrm{GCE}}}{\partial \py} = -\py^{q-1}.
\end{equation}
For $q \in (0,1)$, as $\py \to 0$ (the model is confidently wrong on \emph{another} class),
the gradient magnitude $\py^{q-1} = \py^{-(1-q)} \to +\infty$ but \emph{slower} than CCE
(where the magnitude is $1/\py = \py^{-1}$). For $q=0.7$, the gradient grows as
$\py^{-0.3}$ vs $\py^{-1}$ for CCE — a significant attenuation for small $\py$.

When the label is noisy and the model correctly places low mass on $\tilde y$
(so $p_{\tilde y}$ is small), GCE applies less force than CCE, reducing memorization.

\paragraph{Truncated GCE (TruncGCE).}
Only apply the GCE loss to samples with $p_y < k$ (default $k=0.5$):
\[
  \Lcal_{\mathrm{TrGCE}} = \frac{1-\py^q}{q} \cdot \ind{\py < k}.
\]
Rationale: high-confidence predictions on the observed label are either correct (no gradient needed)
or deeply memorized (should be suppressed, not reinforced).

%------------------------------------------------------------------------------
\subsection{TDPDSCCE — Trimmed DPD Classification Loss}
\label{sec:tdpdscce}
%------------------------------------------------------------------------------

The \textbf{Density Power Divergence} (DPD) for classification is the S-Divergence
with $\lambda = 0$. The full name in our code is \textbf{TDPDSCCE} (Trimmed DPD Sparse CCE).

\paragraph{Definition.}
For parameter $\beta > 0$:
\begin{equation}
  \label{eq:dpd}
  \Lcal_{\mathrm{DPD}}(\pb, y; \beta) = \sum_{k=1}^C p_k^{\beta+1} - \left(1 + \frac{1}{\beta}\right)\py^\beta.
\end{equation}

\paragraph{Derivation from divergence.}
The DPD between discrete distributions $q$ (data) and $p$ (model) is:
\[
  d_\beta(q, p) = \sum_k p_k^{\beta+1} - \frac{\beta+1}{\beta}\sum_k q_k p_k^\beta + \frac{1}{\beta}\sum_k q_k^{\beta+1}.
\]
For the one-hot $q = \eb_y$, $\sum_k q_k p_k^\beta = p_y^\beta$ and the last term is constant,
yielding exactly \cref{eq:dpd} up to the constant.

\paragraph{Gradient analysis.}
\[
  \frac{\partial \Lcal_{\mathrm{DPD}}}{\partial \py} = -\beta \py^{\beta - 1}.
\]
For $\beta < 1$, the gradient magnitude $\beta \py^{\beta-1}$ grows slower than $1/\py$ (CCE)
as $\py \to 0$.

The \textbf{trimmed} version (TDPDSCCE) sorts per-sample losses and keeps only the smallest
$(1-\rho)\cdot n$ of them, discarding the $\rho\%$ most extreme losses (likely noisy samples).

%------------------------------------------------------------------------------
\subsection{Symmetric Cross-Entropy (SCE)}
\label{sec:sce}
%------------------------------------------------------------------------------

Proposed by Wang et al.\ \cite{wang2019sce}:

\paragraph{Definition.}
\begin{equation}
  \label{eq:sce}
  \Lcal_{\mathrm{SCE}}(\pb, y; \alpha, \beta) = \alpha \cdot \Lcal_{\mathrm{CCE}}(\pb, y)
  + \beta \cdot \Lcal_{\mathrm{RCE}}(\pb, y),
\end{equation}
where the \textbf{Reverse Cross-Entropy} treats the one-hot target as the ``model''
and $\pb$ as the ``data'':
\[
  \Lcal_{\mathrm{RCE}}(\pb, y) = -\sum_{k=1}^C p_k \log e_{y,k}
  = -(1-\py)\log\epsilon_0 - C\cdot\log\epsilon_0 + \ldots
\]
In practice, $\eb_y$ is clamped to $\epsilon_0 = 10^{-4}$ so $\log(\eb_{y,\tilde y}) = \log\epsilon_0$
for $\tilde y \neq y$ (a constant), and $\log e_{y,y} = 0$. The effective formula simplifies to:
\[
  \Lcal_{\mathrm{RCE}} \approx -(C-1)\,\py\,\log\epsilon_0 - \ldots \approx \text{const} + C(1-\py).
\]
The code implements:
\begin{equation}
  \Lcal_{\mathrm{SCE}} = \alpha(-\log \py) + \beta \cdot C \cdot (1-\py).
\end{equation}

\paragraph{Why this works.}
The RCE term penalizes distributing probability mass \emph{away} from the correct class —
a complementary signal to CCE. Together, they resist both over-confidence (via CCE)
and under-confidence (via RCE), creating a balanced gradient that is less susceptible
to mislabeled examples.

\paragraph{Scale warning ($\times C$ factor).}
Because of the $C$ multiplier, SCE values on 10-class problems are $\sim 10\times$ larger
than CCE values. \textcolor{red}{\textbf{Never plot SCE and CCE on the same Y-axis!}}

%------------------------------------------------------------------------------
\subsection{S-Divergence Loss (SDIV) — The Institutional Contribution}
\label{sec:sdiv}
%------------------------------------------------------------------------------

SDIV is the core contribution connecting this work to the rSDNet / S-divergence
research programme of Ghosh et al.\ \cite{ghosh2017bern}.

\paragraph{Definition of the S-Divergence.}
For parameters $\beta > 0$, $\lambda \in \R$ satisfying:
\[
  \Asd := 1 + \lambda(1-\beta) > 0, \qquad \Bsd := \beta - \lambda(1-\beta) > 0,
\]
the S-Divergence loss is:
\begin{equation}
  \label{eq:sdiv}
  \boxed{
  \Lcal_{\mathrm{SDIV}}(\pb, y; \beta, \lambda)
  = \frac{1}{\Asd}\sum_{k=1}^C p_k^{\beta+1}
  - \frac{1+\beta}{\Asd \cdot \Bsd}\,\py^{\Bsd}.
  }
\end{equation}
Default in our code (and the rSDNet paper): $\beta=0.05$, $\lambda=-0.8$.

\paragraph{Derivation from the continuous S-Divergence.}
The continuous S-Divergence between distributions $p$ (model) and $q$ (data) is:
\[
  S_{\beta\lambda}(p\|q) = \frac{1}{\Asd}\int p^{\beta+1}\,d\mu
  - \frac{1+\beta}{\Asd \Bsd}\int p^{\Bsd} q^{\lambda(1-\beta)}\,d\mu
  + \text{terms in }q\text{ only}.
\]
For a one-hot $q = \eb_y$, only the mass at $y$ contributes, yielding:
\[
  \int p^{\Bsd} q^{\lambda(1-\beta)}\,d\mu = p_y^{\Bsd} \cdot 1^{\lambda(1-\beta)} = p_y^{\Bsd}.
\]
This recovers \cref{eq:sdiv} (dropping the pure-$q$ constant term).

%------------------------------------------------------------------------------
\subsection{Unification Theorem: S-Divergence Contains All Others}
\label{sec:unification}
%------------------------------------------------------------------------------

\begin{theorem}[S-Divergence Superfamily]
\label{thm:unification}
The SDIV loss in \cref{eq:sdiv} is a superfamily satisfying:
\begin{enumerate}[label=(\alph*)]
  \item \textbf{DPD/TDPDSCCE}: Setting $\lambda = 0$ gives $\Asd = 1$, $\Bsd = \beta$, and
  \[
    \Lcal_{\mathrm{SDIV}}(\pb,y;\beta,0) = \sum_k p_k^{\beta+1} - \frac{\beta+1}{\beta}\,\py^\beta = \Lcal_{\mathrm{DPD}}.
  \]
  \item \textbf{GCE limit}: As $\beta \to 0$ with $\lambda = 0$,
  \[
    \lim_{\beta\to 0}\Lcal_{\mathrm{DPD}} = \lim_{\beta\to 0}\left[\sum_k p_k^{\beta+1} - \frac{\beta+1}{\beta}\py^\beta\right] = -\log\py = \Lcal_{\mathrm{CCE}}.
  \]
  Moreover, for fixed $\beta > 0$ and $\lambda=0$, applying first-order expansion around $\py \approx 1$:
  \[
    \py^\beta \approx 1 + \beta\log\py \quad\Rightarrow\quad \Lcal_{\mathrm{DPD}} \approx -\log\py + O(\beta),
  \]
  confirming the CCE limit.
  \item \textbf{MAE limit}: As $\beta \to 1$ with $\lambda = 0$, $\Bsd = 1$, and
  \[
    \Lcal_{\mathrm{DPD}} = \sum_k p_k^2 - 2\py \to \|p\|_2^2 - 2\py + 1 = \Lcal_{\mathrm{Brier}} - 1,
  \]
  which for the top class is a quadratic analogue of MAE.
\end{enumerate}
\end{theorem}

\begin{proof}
Part (a) follows directly by substituting $\lambda=0$.
Part (b): write $\py^\beta = e^{\beta\log\py}$ and expand around $\beta=0$:
$\py^\beta = 1 + \beta\log\py + O(\beta^2)$.
Then $\sum p_k^{\beta+1} = \sum p_k \cdot p_k^\beta \to \sum p_k = 1$ as $\beta\to 0$.
Substituting: $\Lcal_{\mathrm{DPD}} \to 1 - (1+0)/0 \cdot 1 = \ldots$ (apply L'Hôpital directly):
$\lim_{\beta\to 0}[(1+\beta)/\beta \cdot p_y^\beta] = \lim_{\beta\to 0}[p_y^\beta/\beta + p_y^\beta]$;
the second term gives 1 and the first gives $-\log p_y$ via L'Hôpital.
Part (c) is a straightforward substitution.
\end{proof}

\begin{remark}
This theorem is the \textbf{theoretical justification} for why SDIV is the
primary object of study: a single family recovers CCE-, MAE-, and DPD-type behaviour
through its parameters $(\beta,\lambda)$, enabling a unified analysis of label-noise robustness.
\end{remark}

%==============================================================================
\section{Gradient Landscape and Why These Losses Are Robust}
\label{sec:gradients}
%==============================================================================

The key quantity is $\abs{\partial \Lcal / \partial \py}$: how hard the loss pushes
to increase $p_y$. For a noisy label ($\tilde y$ wrong), a lower gradient means
less memorization.

\begin{table}[h]
\centering\small
\begin{tabular}{@{}llll@{}}
\toprule
\textbf{Loss} & $\partial\Lcal/\partial p_y$ & \textbf{As $p_y \to 0$} & \textbf{Bounded?} \\
\midrule
CCE        & $-1/p_y$              & $\to -\infty$ (fast)   & No  \\
MAE        & $-1$                  & Constant               & Yes \\
GCE$(q)$   & $-p_y^{q-1}$         & $\to -\infty$ (slower) & No  \\
DPD$(\beta)$ & $-\beta\,p_y^{\beta-1}$ & $\to -\infty$ (slower) & No \\
SCE        & $-\alpha/p_y + \beta C$ & $\to -\infty$ but $+\beta C$ resists  & No \\
SDIV       & $-(1+\beta)/A_{\beta\lambda}\cdot p_y^{\Bsd-1}$ & $\to -\infty$ (2-param slower) & No \\
\bottomrule
\end{tabular}
\caption{Gradient at the observed label. Robustness comes from slower divergence
(smaller exponent on $p_y$).}
\label{tab:gradients}
\end{table}

\paragraph{Effective gradient at noisy samples.}
Suppose sample $i$ has noisy label $\tilde y_i \neq y_i^\star$.
After training on clean samples, the model should have small $p_{\tilde y_i}$
(it correctly doubts the noisy label). The gradient forcing the model to memorize this
noise is $\abs{\partial\Lcal/\partial p_{\tilde y_i}}$ evaluated at small $p_{\tilde y_i}$.

From Table~\ref{tab:gradients}: CCE applies the largest force; MAE the smallest;
GCE and SDIV are in between, with the tuning parameters controlling the exact attenuation.

%==============================================================================
\section{Influence Function Analysis}
\label{sec:influence}
%==============================================================================

The \textbf{influence function} (Hampel et al.)\ measures how much the parameter
estimate $\hat\theta$ would change if training point $(\xb_i, y_i)$ were up-weighted:
\[
  \mathrm{IF}(\xb_i, y_i; F_n) = H^{-1} \nabla_\theta \ell(\pb(\xb_i;\hat\theta), y_i),
\]
where $H = \nabla_\theta^2 \E[\ell]$ is the Hessian. Smaller gradient norm $\Rightarrow$
smaller influence $\Rightarrow$ less memorization.

\begin{proposition}
\label{prop:influence}
For the GCE loss with parameter $q \in (0,1)$, the influence of a maximally noisy sample
(with $p_{\tilde y}(\hat\theta) = 1/(C-1+\epsilon)$ after partial training) satisfies:
\[
  \norm{\nabla_\theta \Lcal_{\mathrm{GCE}}} / \norm{\nabla_\theta \Lcal_{\mathrm{CCE}}}
  \leq \left(\frac{1}{C-1}\right)^{q-1} \xrightarrow[C\to\infty]{} 0^+.
\]
Thus in large-class problems, GCE suppresses the influence of uniform-distribution
noisy samples by a factor of $(C-1)^{1-q}$.
\end{proposition}

%==============================================================================
\section{Forward Correction (Label-Correction Losses)}
\label{sec:forward}
%==============================================================================

With known or estimated transition matrix $\Tb$, Patrini et al.\ \cite{patrini2017} derive:

\paragraph{Corrected predictive distribution.}
Given clean predicted distribution $\pb$ and noise model $\Tb$:
\[
  \pb' = \Tb^\top \pb, \qquad p'_{\tilde y} = \sum_j T_{j\tilde y} p_j.
\]
The \textbf{forward-corrected loss} is:
\begin{equation}
  \label{eq:forwardT}
  \Lcal_{\mathrm{Fwd}}(\pb, \tilde y; \Tb) = -\log \big[\Tb^\top \pb\big]_{\tilde y}.
\end{equation}

\paragraph{Estimation of $\hat{\Tb}$ (anchor-point heuristic).}
Patrini et al.\ \cite{patrini2017} show that for each class $c$, the sample with
highest predicted probability $p_c$ is an ``anchor point'' — likely a clean example.
The $c$-th row of $\hat{\Tb}$ is estimated as the softmax output at this anchor:
\[
  \hat T_{c,\cdot} = \pb\big(\xb^{(c)}_{\mathrm{anchor}}\big), \quad
  \xb^{(c)}_{\mathrm{anchor}} = \arg\max_{\xb \in \mathcal{D}} p_c(\xb;\hat\theta).
\]

%==============================================================================
\section{Dual Robustness: Label Noise and Adversarial Perturbation}
\label{sec:dual}
%==============================================================================

\begin{tcolorbox}[colback=green!5!white,colframe=green!60!black,title=Novel Contribution]
This section is a novel theoretical result not present in the existing codebase
or write-up. It provides the mathematical foundation for simultaneously evaluating
robustness against both label noise and adversarial perturbations.
\end{tcolorbox}

Let $\ell : \R^C \times \mathcal{Y} \to \R_+$ be a classification loss.
Define the \textbf{noisy risk} and \textbf{adversarial risk} respectively as:
\[
  R_\eta(\theta) = \E_{(\xb,\tilde y)\sim P_\eta}\!\left[\ell\big(\pb(\xb;\theta),\tilde y\big)\right],
\]
\[
  R_\epsilon(\theta) = \E_{(\xb,y)\sim P}\!\left[\max_{\norm{\boldsymbol{\delta}}_\infty \leq \epsilon}
  \ell\big(\pb(\xb+\boldsymbol{\delta};\theta),y\big)\right].
\]

\begin{theorem}[Dual Robustness Bound — Novel]
\label{thm:dual}
Let $\ell$ be a loss satisfying the \textbf{$\rho$-gradient bound}:
$\abs{\partial\ell/\partial p_y} \leq \rho$ for all $(p_y \in (0,1])$.
Then simultaneously:
\begin{align}
  R_\eta(\theta) &\leq R_0(\theta) + \eta \cdot \rho \cdot \|\pb\|_\infty, \label{eq:noise-bound}\\
  R_\epsilon(\theta) &\leq R_0(\theta) + \epsilon\sqrt{d}\cdot L_f \cdot \rho, \label{eq:adv-bound}
\end{align}
where $R_0(\theta) = \E[\ell(\pb(\xb),y^\star)]$ is the clean risk, $L_f$ is the
Lipschitz constant of $f_\theta$, and $d$ is the input dimension.
\end{theorem}

\begin{proof}[Proof sketch]
For \cref{eq:noise-bound}: Write $R_\eta = (1-\eta)R_0 + \eta R_\perp$ where
$R_\perp = \E[\ell(\pb(\xb),y_\perp)]$ for uniform wrong class $y_\perp$.
By the gradient bound, $\ell(\pb,y_\perp) \leq \ell(\pb,y) + \rho|p_{y_\perp} - p_y| \leq R_0 + \rho$.
So $R_\eta \leq R_0 + \eta\rho$.

For \cref{eq:adv-bound}: By Lipschitz of the composition, perturbing input by $\boldsymbol{\delta}$
changes logits by at most $L_f\norm{\boldsymbol{\delta}}$. The loss changes by at most
$\rho \cdot \Delta p_y$ where $\Delta p_y \leq L_\softmax \cdot L_f \norm{\boldsymbol{\delta}}_2
\leq L_f\sqrt{d}\epsilon$ (for $\|\boldsymbol{\delta}\|_\infty \leq \epsilon$).
\end{proof}

\begin{corollary}
MAE (with $\rho=1$) achieves the \textbf{tightest dual robustness bound}. CCE
(with $\rho = 1/p_y^{\min} \gg 1$) achieves the worst. GCE and SDIV are between:
$\rho \approx p_y^{\min\,(q-1)}$ and $\rho \approx p_y^{\min\,(\Bsd-1)}$ respectively.
\end{corollary}

This justifies our experimental design: we test both label noise (Battery B) and
adversarial attacks (Battery C) with all losses, and hypothesize that losses with
smaller effective $\rho$ dominate both robustness metrics.

%==============================================================================
\section{Loss Function Scales: The Plotting Problem and Its Solution}
\label{sec:scales}
%==============================================================================

\textcolor{red}{\textbf{Critical issue with previous plots}}: the Y-axis in training-loss
curves compared losses with incompatible scales. Here are the correct ranges:

\begin{table}[h]
\centering\small
\begin{tabular}{@{}lllp{5cm}@{}}
\toprule
\textbf{Loss} & \textbf{Range} & \textbf{At init ($p_y=1/C$)} & \textbf{Correct Y-axis label} \\
\midrule
CCE    & $[0, +\infty)$ & $\log C \approx 2.3$ ($C=10$)         & ``Cross-Entropy (nats)'' \\
MAE    & $[0, 1]$       & $(C-1)/C \approx 0.9$                 & ``$1 - p_y$'' \\
GCE$(q)$ & $[0, 1/q]$   & $(1-C^{-q})/q$                       & ``$(1-p_y^q)/q$'' \\
TDPDSCCE & $\R$ (can be negative) & $C^{-\beta}(1 - (1+1/\beta))$  & ``DPD loss (unitless)'' \\
SCE    & $[0, +\infty)$ but $\times C$ & $\alpha\log C + \beta C (1-1/C)$ & ``SCE ($\alpha$ CCE + $\beta$ RCE)'' \\
SDIV   & Could be negative & depends on $(\beta,\lambda)$ & ``S-Div loss ($\beta$=0.05, $\lambda$=-0.8)'' \\
\bottomrule
\end{tabular}
\caption{Loss function scales. \textbf{Do not compare on a common y-axis}.
Each should have its own scale, or a normalized $[0,1]$ min-max scale for comparison.}
\label{tab:scales}
\end{table}

\paragraph{Correct plotting strategy.}
\begin{enumerate}
  \item \textbf{Training loss curves}: plot each loss family on a \emph{separate subplot}
  with an accurately labeled Y-axis (see Table \ref{tab:scales}).
  \item \textbf{Accuracy curves}: all losses can share one plot (accuracy is comparable).
  \item \textbf{Normalized comparison}: min-max normalize each loss to $[0,1]$ within
  its own training run if you must show all curves together. Always state in the caption
  that values are normalized.
  \item \textbf{Confusion matrices}: always use \texttt{normalize='true'} (row normalization),
  which gives \emph{recall per class} as values in $[0,1]$. Raw counts are meaningless
  without comparing to class sizes.
\end{enumerate}

%==============================================================================
\section{Experimental Design for Conference Submission}
\label{sec:expdesign}
%==============================================================================

\subsection{Battery Structure}

\begin{table}[h]
\centering\small
\begin{tabular}{@{}lllll@{}}
\toprule
\textbf{Battery} & \textbf{Name} & \textbf{Variable} & \textbf{Metric} & \textbf{Claim to test} \\
\midrule
A & Clean baseline   & -- & Accuracy, $F_1$ macro & All losses competitive at $\eta=0$ \\
B & Uniform noise    & $\eta \in \{0,0.2,0.4,0.6\}$ & Test acc & SDIV/GCE $>$ CCE at high $\eta$ \\
C & Adversarial FGSM & $\epsilon/255 \in \{0,1,2,4,8\}$ & Acc after FGSM & Bounded-grad losses more robust \\
D & $q$-sweep (GCE)  & $q \in \{0, 0.4, 0.7, 1.0\}$ & Acc, val loss & Optimal $q$ depends on $\eta$ \\
E & $(\beta,\lambda)$ surface & $\beta \in \{0.02,...,0.5\}$, $\lambda$ & 3D acc surface & Optimal SDIV hyperparams \\
\bottomrule
\end{tabular}
\caption{Experimental batteries and their claims.}
\end{table}

\subsection{Datasets in Our Pipeline}

The three dataset families in our code map to distinct research claims:

\begin{itemize}
  \item \textbf{MNIST / CIFAR-10} (Part 2, TF Keras): \emph{controlled benchmarks} to verify
  known theory results from the literature (GCE $>$ CCE at $\eta > 0.2$, etc.)
  \item \textbf{PathMNIST / DermaMNIST + CLIP} (Part 4, PyTorch zero-shot):
  \emph{domain-shift analysis} under CLIP zero-shot logits. Novel because no fine-tuning
  means the loss comparison is purely metric-level.
  \item \textbf{Emotion / HealthFact / PubMedQA + BERT} (Part 3, PyTorch fine-tuning):
  \emph{NLP under label noise}, a significantly under-explored regime compared to vision.
\end{itemize}

\subsection{Why You Must Train from Scratch for Each Loss}

\begin{tcolorbox}[colback=yellow!5!white,colframe=orange!80!black,title=Key Question Answered]
\textbf{Q: Can I use pretrained weights and just switch the loss?}\\[4pt]
\textbf{A (short): Mostly no.}\\[4pt]
The pretrained weights $\theta_0$ were obtained by minimizing CCE. They encode a specific
\emph{loss-specific optimization trajectory}: the gradient flow under CCE shaped the entire
representation. Switching to GCE or SDIV at this pretrained point and doing a \emph{brief}
fine-tune will measure optimization dynamics, not the true converged representation each
loss learns.\\[4pt]
\textbf{Exception — CLIP-based zero-shot (Part 4)}: the logits are \emph{fixed} (no training
at all). Here you compute the loss as a metric on frozen zero-shot outputs, which is valid
and novel.\\[4pt]
\textbf{Exception — BERT fine-tuning (Part 3)}: you fine-tune from a pretrained checkpoint
for \emph{each loss separately}. This is standard practice: all BERT papers start from
\texttt{bert-base-uncased} or domain-specific checkpoints — the loss determines the
fine-tuning trajectory. This is valid because fine-tuning is long enough to escape the
CCE-shaped basin and converge under the new loss.
\end{tcolorbox}

\subsection{Novel Research Directions for Core A* Publication}

The individual losses are established. The \textbf{novel contributions} for a spotlight-award
paper are:

\begin{enumerate}
  \item \textbf{Unified S-Divergence Framework} (Theorem~\ref{thm:unification}):
  First systematic comparison across the $(\beta,\lambda)$ surface for vision, NLP, and
  multimodal models under the same experimental protocol.

  \item \textbf{Dual Robustness Benchmark} (Theorem~\ref{thm:dual}):
  Simultaneous evaluation against label noise AND adversarial perturbation. Existing
  literature treats these separately. We show SDIV provides a principled tradeoff.

  \item \textbf{CLIP Zero-Shot Loss Evaluation}:
  Using CLIP as a \emph{fixed feature extractor} and studying how loss-function metrics
  change under label noise is novel: no prior work evaluates robust losses in the
  zero-shot setting for medical imaging.

  \item \textbf{Cross-Modal Robustness Comparison}:
  Same loss family (SDIV, GCE) evaluated on MNIST (visual), PubMedQA (biomedical NLP),
  and PathMNIST (medical vision) within a single paper — showing generalizability of
  the institutional S-divergence framework.

  \item \textbf{Curriculum Loss Annealing} (proposed):
  Schedule the loss parameter as a function of training progress:
  $q(t) = q_{\max} \cdot (1 - t/T)$ — start with high $q$ (robust MAE-like) and anneal
  toward low $q$ (efficient CCE-like). This mirrors known curriculum effects and has
  theoretical justification from \cref{thm:dual}.
\end{enumerate}

%==============================================================================
\section{Mathematical Appendix: Key Proofs and Derivations}
\label{sec:appendix}
%==============================================================================

\subsection{Softmax Jacobian}

\[
  \frac{\partial p_j}{\partial z_k} = p_j(\delta_{jk} - p_k),
\]
where $\delta_{jk}$ is the Kronecker delta.

Using this, the CCE gradient w.r.t.\ logits becomes:
\[
  \frac{\partial \Lcal_{\mathrm{CCE}}}{\partial z_k} = \sum_j \frac{\partial \Lcal}{\partial p_j}\frac{\partial p_j}{\partial z_k}
  = -\frac{1}{\py}\sum_j \delta_{jy}\frac{\partial p_j}{\partial z_k}
  = -\frac{1}{\py}\cdot \py(\delta_{yk} - p_k)
  = p_k - \delta_{yk},
\]
recovering the familiar result.

\subsection{GCE Gradient w.r.t.\ Logits}

\begin{align}
  \frac{\partial \Lcal_{\mathrm{GCE}}}{\partial z_k}
  &= \frac{\partial}{\partial z_k}\frac{1-\py^q}{q}
  = -p_y^{q-1}\frac{\partial p_y}{\partial z_k}
  = -p_y^{q-1}\cdot p_y(\delta_{yk} - p_k) \\
  &= \begin{cases}
  -p_y^q(1-p_y) & k=y \\
  p_y^q\cdot p_k & k\neq y.
  \end{cases}
\end{align}

Compare with CCE: the factor $p_y^q$ suppresses the gradient when $p_y$ is small
(low confidence on the label). For noisy labels where the correct answer has been
forgotten, $p_y$ is small so the gradient is suppressed --- reducing memorization.

\subsection{SDIV Gradient w.r.t.\ Logits}

Let $g_k = p_k^{\beta+1}$ and $h = p_y^{\Bsd}$. Then:
\begin{align}
  \frac{\partial \Lcal_{\mathrm{SDIV}}}{\partial z_k}
  &= \frac{1}{\Asd}(\beta+1)p_k p_k^{\beta-1}(p_k - p_k^2/p_j) - \frac{1+\beta}{\Asd\Bsd}\Bsd p_y^{\Bsd-1}\frac{\partial p_y}{\partial z_k}
\end{align}

Simplifying using the softmax Jacobian:
\begin{equation}
  \frac{\partial \Lcal_{\mathrm{SDIV}}}{\partial z_k}
  = \frac{1+\beta}{\Asd}\left[p_k^\beta(p_k - p_k^2) - p_y^{\Bsd-1}p_y(\delta_{yk}-p_k)/1\right].
\end{equation}

The suppression factor is $p_y^{\Bsd-1}$ — with $\Bsd = \beta - \lambda(1-\beta) = 0.05 - (-0.8)(0.95) = 0.81$
for our default parameters. So the suppression goes as $p_y^{-0.19}$, much weaker than
CCE's $p_y^{-1}$.

%==============================================================================
\section{Research Calendar and Next Steps}
\label{sec:calendar}
%==============================================================================

\begin{table}[h]
\centering\small
\begin{tabular}{@{}llll@{}}
\toprule
\textbf{Date} & \textbf{Task} & \textbf{Deliverable} & \textbf{Status} \\
\midrule
30 March 2026 & MNIST/CIFAR experiments & CSV + PNG plots & Done \\
12 April 2026 & Theory write-up; code refactor & This .tex; new .py/.ipynb & In progress \\
19 April 2026 & RunPod full Part3+Part4 runs & Results CSV (5 seeds) & Next \\
26 April 2026 & Dual-robustness experiments & Battery B+C joint plots & \\
3 May 2026    & Curriculum annealing pilot & Loss schedule ablation & \\
10 May 2026   & First draft of paper & NeurIPS/ICML format & \\
\bottomrule
\end{tabular}
\end{table}

%==============================================================================
\begin{thebibliography}{99}

\bibitem{zhang2018gce}
Z.~Zhang and M.~R. Sabuncu.
\newblock Generalized Cross Entropy Loss for Training Deep Neural Networks with Noisy Labels.
\newblock In \emph{NeurIPS}, 2018.

\bibitem{wang2019sce}
Y.~Wang et al.
\newblock Symmetric Cross Entropy for Robust Learning with Noisy Labels.
\newblock In \emph{ICCV}, 2019.

\bibitem{patrini2017}
G.~Patrini et al.
\newblock Making Deep Neural Networks Robust to Label Noise: A Loss Correction Approach.
\newblock In \emph{CVPR}, 2017.

\bibitem{ghosh2017bern}
A.~Ghosh et al.
\newblock A Generalized Divergence for Statistical Inference.
\newblock \emph{Bernoulli}, 23(4A):2746--2783, 2017.

\bibitem{ghosh2026arxiv}
A.~Ghosh and S.~Jana.
\newblock Provably robust learning of regression neural networks using $\beta$-divergences.
\newblock arXiv:2602.08933, 2026.

\bibitem{yang2023medmnist}
J.~Yang et al.
\newblock MedMNIST v2 --- A large-scale lightweight benchmark for 2D and 3D biomedical image classification.
\newblock \emph{Scientific Data}, 10(1):41, 2023.

\bibitem{goodfellow2014}
I.~J. Goodfellow, J.~Shlens, and C.~Szegedy.
\newblock Explaining and Harnessing Adversarial Examples.
\newblock In \emph{ICLR}, 2015.

\bibitem{han2018coteaching}
B.~Han et al.
\newblock Co-teaching: Robust training of deep neural networks with extremely noisy labels.
\newblock In \emph{NeurIPS}, 2018.

\end{thebibliography}

\end{document}
